% a06_pgate_c: Step 3(F), fifth part, end (input from sec/a06_BALtorus_ct8.tex): the final step, the field-dependent part of
% Ba2 (2.116), as Theorem thm:a06f:C2 at k = K with an empty inner domain.

\emph{The final step.} Let $K$ be the last step, $l_K=t_K=L^K\eps\in(\eps_0/L,\eps_0]$ (\cref{thm:a06:F}), in scale-$K$ units, with the Setting of \cref{thm:a06f:C2} at $k=K$: $\Omega_A=B^K(\Lambda^{(K-1)'}_2)$, $\Omega_1=B^K(\Lambda^{(K-1)'}_7)$, $\tilde B=A^{(K)}=\theta_KA^{(K)}$ on $\Omega_A$, $\Phi=\phi^{(K)}$ of \cite[(2.56)]{Ba2} at $k=K$, under \cite[(2.55)]{Ba2} at step $K$. No set $\Lambda^{(K)}_i$ is defined at the final step: \cite[(2.116)]{Ba2} bounds $P^{(K)}$ on its whole domain $\Omega_1$, and its right side is the volume $|\Lambda^{(K-1)'}_7|$ of that domain, as in the last term of \cite[(2.117)]{Ba2}. The chain is run with $E_0+E_K$ (\cref{lem:a06:CF}), so no step polynomial contains $\Efin$, which is the global factor bounded by \cref{thm:a06:F}.

\begin{proposition}[{the final step, \cite[(2.116)]{Ba2}}]\label{prop:a06g:FS}
For every $\beta>0$,
\[
P^{(K)}(\Omega_1,\tilde B,\varphi)=-\lambda_K\sum_{x\in\Omega_1}\eta^3|\Phi(x)|^4+R_{\mathrm{field}}+c_K(\Omega_1),\qquad|R_{\mathrm{field}}|\le C_\beta t_K^{1/2-\beta}|\Lambda^{(K-1)'}_7|,
\]
with $c_K(\Omega_1)$ the field-independent part, $|c_K(\Omega_1)|\le Ct_K(1+\ln t_K^{-1})^5|\Lambda^{(K-1)'}_7|\le C_\beta t_K^{1-\beta}|\Lambda^{(K-1)'}_7|$. Hence $P^{(K)}(\Omega_1,\tilde B,\varphi)\le C_\beta t_K^{1/2-\beta}|\Lambda^{(K-1)'}_7|$, which is \cite[(2.116)]{Ba2} with $\kappa_0=\frac12-\beta$. The constants are those of \cref{thm:a06f:C2}, free of $\eps$, $K$, $\eps_0$, the torus and the history.
\end{proposition}

\begin{proof}
Run the proof of \cref{thm:a06f:C2} at $k=K$ with the inner domain $\Omega'_1\defeq\emptyset$.

\emph{Step 1: the empty domain.} By \cref{lem:a06f:C21} the vertex domain enters only through $\mathbf 1_{\Omega'_1}$ on the vertices \cite[(1.6)--(1.7)]{Ba3}, through the gate on the $A'$-legs ($\supp g\subset\Omega'_1$) and through the localised counterterms. With $\Omega'_1=\emptyset$ every vertex vanishes and every $A'$-leg with it, so $P^{(K)}(\emptyset,\tilde B,\varphi)$ is field-independent and equals $c_K(\emptyset)$, which is $0$: on the chain run with $E_0+E_K$, $c_k(\Omega')=-E_{K,k}+a_k(\Omega')$ (\cref{lem:a06e:VK0}), $E_{K,K}=E_K-E_K=0$, and $a_K(\emptyset)$ has no pair vertex to place. In \cref{lem:a06f:C22}, $D_{\mathrm{good}}=\emptyset$ and $\widetilde{\mathrm{Col}}=\Omega_1$: every localised family is marked, and $P^{(K)}(\Omega_1)=P^{(K)}(\Omega_1)-P^{(K)}(\emptyset)$ is the sum over all localised families, rooted at a cube of $\Omega_1$ ($|\Omega_1|=|\Lambda^{(K-1)'}_7|$ choices).

\emph{Step 2: the three kinds of family.} The families with $W\ge2$ are bounded by \cref{lem:a06f:C25} at $k=K$: Proposition~1 on the lines $G_K(\Omega_A,\tilde B)$ with the external scalar fields of \cref{lem:a06f:C23} at $k=K$ (which uses \cite[(2.55)]{Ba2} at step $K$ and \cite[(2.24)--(2.25)]{Ba1} at points of $\Omega_1$, at depth $\ge r(t_K)-2\ge R_0$ in $\Omega_A$, $\tilde B$ regular by \cite[(2.60)]{Ba2}), the \cite[(1.7)]{Ba3} vertices through \cref{thm:a06e:T0} and \cref{lem:a06e:T1}, the sum by \cref{lem:a06d:V11} with the root in $\Omega_1$: in all $\le C_\beta t_K^{1/2-\beta}|\Lambda^{(K-1)'}_7|$. The families with $W=0$ are \cref{lem:a06f:C26}: the lone \cite[(1.6)]{Ba3} vertex, whose sum over the vertex domain is sharp, gives exactly $-\lambda_K\sum_{x\in\Omega_1}\eta^3|\Phi(x)|^4$; the single-line terms give $\le C_{\mathrm{SL}}e^2\lambda^{-1/2}p^4t_K^{1/2}$ per cube of $\Omega_1$ by \cite[Lemma~2.4, (2.65)--(2.66)]{Ba2} at every $y\in\Lambda^{(K-1)'}_7$ (\cref{thm:a06:BALA}(c)) and the slice kernels of \cref{lem:a06e:SL}, in all $\le Ct_K^{1/2-\beta}|\Lambda^{(K-1)'}_7|$. The field-independent families are \cref{lem:a06f:C27}(ii) with $\Omega'_1=\emptyset$: the vacuum families with a cube meeting $\Omega_1$ have $W\ge4$, and $|c_K(\Omega_1)|=|c_K(\Omega_1)-c_K(\emptyset)|\le Ct_K(1+\ln t_K^{-1})^5|\Lambda^{(K-1)'}_7|$; at zero background this is \cref{cor:a06d:V13}, and on the chain of \cref{lem:a06:CF} it is the bound of \cref{thm:a06e:VK}(a) on $a_K(\Omega_1)$, whose proof does not use $k<K$.

\emph{Step 3: the inputs at $k=K$.} Every input of \cref{thm:a06f:C2} is available at the final step: \cite[(2.55)]{Ba2} holds at every step, $K$ included; Proposition~1's $O(1)$ does not depend on the step; the layer distances used are those of the step-$(K-1)$ sets, $\dist(\Omega_1,\Omega_A^c)\ge r(t_K)-2$; and no set $\Lambda^{(K)}_i$ enters, since the inner domain is empty. The one-sided bound follows by dropping the quartic term, which is $\le0$.
\end{proof}

The same bound, with $|\Lambda^{(K-1)'}_7|$ replaced by $L^3|T^{(K)}|\ge|\Lambda^{(K-1)'}_7|$, is the $k=K$ term of the exponents of \cite[(3.42)]{Ba2} as \cref{prop:a06f:RS} carries it. This completes the field-dependent part of the second sentence of \cite[p.~579]{Ba2} (\cref{prop:a06g:HP}) and of the final step (\cref{prop:a06g:FS}); with \cref{thm:a06f:C2} (the collar), the three transfers of the mechanism of \cite[Prop.~2.1]{Ba2} that \cite[p.~570]{Ba2} announces (``similar conclusions hold for other expressions which will be included in the action in the later stages of the procedure'') are written.
